\documentclass[11pt,a4paper]{article}
\usepackage[utf8]{inputenc}
\usepackage[spanish]{babel}
\usepackage{listings}
\usepackage{xcolor}
\usepackage{hyperref}
\usepackage{geometry}
\geometry{margin=1in}

% Configuración de colores para los fragmentos de código
\definecolor{codegreen}{rgb}{0,0.6,0}
\definecolor{codegray}{rgb}{0.5,0.5,0.5}
\definecolor{codepurple}{rgb}{0.58,0,0.82}
\definecolor{backcolour}{rgb}{0.08,0.11,0.17}
\definecolor{codefg}{rgb}{0.9,0.9,0.9}

\lstdefinestyle{mystyle}{
    backgroundcolor=\color{backcolour},   
    commentstyle=\color{codegreen},
    keywordstyle=\color{orange},
    numberstyle=\tiny\color{codegray},
    stringstyle=\color{codepurple},
    basicstyle=\ttfamily\footnotesize\color{codefg},
    breakatwhitespace=false,         
    breaklines=true,                 
    captionpos=b,                    
    keepspaces=true,                 
    numbers=left,                    
    numbersep=5pt,                  
    showspaces=false,                
    showstringspaces=false,
    showtabs=false,                  
    tabsize=2
}

\lstset{style=mystyle}

\title{Documentación de Actualización Estacional y Manipulación del DOM}
\author{AEGIS Wiki - Elden Ring}
\date{\today}

\begin{document}

\maketitle

\section{Introducción}
Esta documentación detalla los desarrollos técnicos y las características implementadas para la actualización de temporada de la \textbf{Wiki de Elden Ring (AEGIS Wiki)}, enfocado en eventos del puntero, tematización del calendario y administración interactiva del DOM.

\section{Evento de Destacado mediante Puntero (Spotlight Glow)}

\subsection{Descripción}
Para la \textbf{Temporada de la Gracia Dorada (Verano)}, se implementó un evento interactivo que detecta la posición del cursor sobre las tarjetas de la interfaz para proyectar un foco o resplandor de luz dorado (\textit{spotlight}) que sigue el movimiento del puntero, acompañado de una transición de escala y cambio de color del borde.

\subsection{Detalles Técnicos e Implementación}
El evento se programó en Angular mediante una directiva \textit{standalone} (\texttt{HighlightDirective}), utilizando escuchadores de eventos nativos de JavaScript para dispositivos señaladores (\textit{Pointer Events}):

\begin{enumerate}
    \item \textbf{Directiva (\texttt{highlight.directive.ts})}:
    \begin{itemize}
        \item \textbf{\texttt{pointerenter}}: Activa la clase \texttt{.seasonal-highlight-active} que inicia la transición visual (escala, borde dorado y opacidad del brillo).
        \item \textbf{\texttt{pointermove}}: Registra las coordenadas relativas \texttt{clientX} y \texttt{clientY} respecto al contenedor usando \texttt{getBoundingClientRect()}, y actualiza dos variables de entorno CSS (\texttt{--mouse-x} y \texttt{--mouse-y}).
        \item \textbf{\texttt{pointerleave}}: Remueve la clase activa y oculta el resplandor para optimizar el rendimiento y limpiar el DOM visual.
    \end{itemize}
\end{enumerate}

\begin{lstlisting}[language=Java, caption=Fragmento de la directiva highlight.directive.ts en TypeScript]
// Fragmento de highlight.directive.ts
@HostListener('pointermove', ['$event'])
onPointerMove(event: PointerEvent) {
  const rect = this.el.nativeElement.getBoundingClientRect();
  const x = event.clientX - rect.left;
  const y = event.clientY - rect.top;
  this.renderer.setStyle(this.el.nativeElement, '--mouse-x', `${x}px`);
  this.renderer.setStyle(this.el.nativeElement, '--mouse-y', `${y}px`);
}
\end{lstlisting}

\begin{enumerate}
    \item[2.] \textbf{Estilos CSS (\texttt{styles.css})}:
    Utiliza un gradiente radial flotante posicionado con las variables actualizadas por el evento JavaScript:
\end{enumerate}

\begin{lstlisting}[language=HTML, caption=Fragmento de estilos CSS para el spotlight]
.seasonal-highlight::before {
	content: '';
	position: absolute;
	top: 0; left: 0; right: 0; bottom: 0;
	background: radial-gradient(
		80px circle at var(--mouse-x, 0px) var(--mouse-y, 0px),
		rgba(202, 165, 81, 0.15),
		transparent 80%
	);
	opacity: 0;
	transition: opacity 0.3s ease;
	pointer-events: none;
}
.seasonal-highlight-active::before {
	opacity: 1;
}
\end{lstlisting}

\section{Menú Desplegable Basado en Eventos de Puntero (Dropdown)}

\subsection{Descripción}
Un menú de navegación en la barra superior (``Wiki Categorías'') que se despliega dinámicamente al pasar el puntero sobre él y se oculta cuando el puntero sale de su área delimitada. Esto permite una navegación rápida y fluida sin requerir clics innecesarios.

\subsection{Detalles Técnicos e Implementación}
\begin{itemize}
    \item El contenedor del menú se enlazó a los eventos JavaScript/Angular \texttt{(pointerenter)} y \texttt{(pointerleave)} en la plantilla \texttt{app.html}.
    \item Dichos eventos modifican el estado de un Signal reactivo (\texttt{dropdownVisible}).
    \item Al cambiar el estado a \texttt{true}, el DOM renderiza condicionalmente el contenedor del menú desplegable (\texttt{.nav-dropdown-menu}), el cual posee un diseño \textit{glassmorphic} y una animación de entrada de desvanecimiento y deslizamiento (\texttt{slideDownFade}).
    \item Al hacer clic en un enlace de categoría o retirar el puntero, se dispara el evento para cerrar el menú inmediatamente.
\end{itemize}

\section{Tematización de Interfaz basada en el Calendario}

\subsection{Descripción}
El sitio web cuenta con un detector de fecha calendarizado que adapta la cabecera, el banner y el esquema de colores de la interfaz dependiendo del mes actual en el sistema del cliente.

\subsection{Calendario de Eventos}
\begin{itemize}
    \item \textbf{Verano (Meses Junio, Julio, Agosto)}: \textit{Temporada de la Gracia Dorada (Golden Grace Season)}.
    \begin{itemize}
        \item Estilo: Resplandor dorado, ámbar y azul cobalto oscuro.
        \item Banner: ``AEGIS Wiki $\cdot$ Solsticio de Oro - Temporada de la Gracia Dorada''.
    \end{itemize}
    \item \textbf{Otoño (Meses Octubre, Noviembre)}: \textit{Temporada del Fuego de la Zarza (Shadow Bramble Season)}.
    \begin{itemize}
        \item Estilo: Resplandor carmesí, ceniza y negro carbón.
        \item Banner: ``AEGIS Wiki $\cdot$ Cosecha de Sombras - Temporada del Fuego de la Zarza''.
    \end{itemize}
    \item \textbf{Invierno (Meses Diciembre, Enero, Febrero)}: \textit{Temporada de la Helada del Gigante (Frost Glintstone Season)}.
    \begin{itemize}
        \item Estilo: Resplandor azul glacial, escarcha y plateado.
        \item Banner: ``AEGIS Wiki $\cdot$ Invierno Helado - Temporada de la Helada del Gigante''.
    \end{itemize}
    \item \textbf{Resto del año}: Tema por defecto del sitio.
\end{itemize}

\subsection{Flujo de Ejecución}
En el ciclo de inicio \texttt{ngOnInit()} de \texttt{app.ts}, se evalúa:
\begin{lstlisting}[language=Java]
const currentMonth = new Date().getMonth();
// Asigna clases como 'theme-golden-grace' y modifica los signals del banner
\end{lstlisting}
La clase resultante se acopla al tag principal del DOM \texttt{<main [class]="currentTheme()">}, heredando los colores correspondientes de las variables CSS.

\section{Administrador de Tareas: Bitácora de Misiones del Sinluz}

\subsection{Descripción}
Para satisfacer la manipulación avanzada de la estructura, contenido y estilos del DOM, se implementó una página interactiva de administración de misiones a la cual se accede desde la barra lateral o el menú desplegable.

\subsection{Funciones Implementadas del DOM y Eventos}
\begin{enumerate}
    \item \textbf{Creación}: El usuario llena un formulario. El botón de envío se valida activamente. Al presionar ``Añadir Misión'', se crea dinámicamente una nueva tarjeta en la rejilla del DOM.
    \item \textbf{Validaciones}: Validación en tiempo real del título (mínimo 4 letras) y el objetivo (mínimo 10 letras), mostrando mensajes de error dinámicos en el DOM y deshabilitando el botón de envío si no se cumplen.
    \item \textbf{Modificación (Edición)}: Al presionar ``Editar'', se cargan los datos en el formulario para modificarlos. Al guardar, se actualiza el contenido de la tarjeta en el DOM.
    \item \textbf{Actualización de Estados (Completado)}: Al presionar ``Completar/Reabrir'', se altera el estilo de la tarjeta (se tacha el título, se cambia la opacidad y el color del chip de estado) y se actualizan dinámicamente las estadísticas del progreso.
    \item \textbf{Eliminación}: El botón ``Eliminar'' remueve físicamente el nodo del DOM y recalcula las estadísticas globales de misiones de inmediato.
    \item \textbf{Estadísticas Dinámicas}: Un panel circular calcula de forma computada el porcentaje de misiones completas y actualiza los contadores en tiempo real.
\end{enumerate}

\end{document}
